\documentclass[12pt]{article}

\usepackage{sbc-template}
\usepackage{graphicx,url}
\usepackage[utf8]{inputenc}
\usepackage[brazil]{babel}
\usepackage{pgfplots}
\usepackage{pgfplotstable}
\pgfplotsset{compat=1.18}

\sloppy

\title{Estudo Experimental de Algoritmos de Ordenacao em C: analise de comparacoes, trocas e tempo de execucao}

\author{Fabricio Rocha\inst{1}}

\address{Trabalho Pratico -- Estruturas de Dados II\\
  \email{fabricio@example.com}
}

\begin{document}

\maketitle

\begin{abstract}
This paper presents a revised empirical evaluation of 13 sorting algorithms implemented
in C, following the updated assignment requirements. Experiments were performed for
instances with 10,000 and 1,000,000 elements under three input patterns: random,
ascending and descending. For each execution, we report comparisons, swaps and runtime.
The study includes the two required Quicksort variants (center pivot and median-of-three
pivot), and adopts a maximum execution time of 24 hours per algorithm-instance pair,
marking non-finished runs as timeout.
\end{abstract}

\begin{resumo}
Este artigo apresenta uma versao revisada da avaliacao experimental de 13 algoritmos de
ordenacao implementados em C, em conformidade com o enunciado atualizado. Os
algoritmos foram executados sobre instancias com 10.000 e 1.000.000 de elementos em
tres padroes de entrada: aleatoria, ordenada crescente e ordenada decrescente. Em cada
execucao foram coletados numero de comparacoes, numero de trocas e tempo de execucao.
O estudo contempla as duas variantes exigidas de Quicksort (pivo central e mediana de
tres) e adota tempo maximo de 24 horas por par algoritmo-instancia, registrando como
timeout as execucoes nao finalizadas nesse limite.
\end{resumo}

\section{Introducao}
O objetivo deste trabalho foi comparar, de forma controlada, diferentes estrategias de
ordenacao estudadas na disciplina. Foram analisados os algoritmos: bolha,
bolha com parada, insercao, insercao binaria, insercao ternaria, shell,
selecao, heap, quicksort com pivo central, quicksort com mediana de tres,
merge, radix e bucket.

O foco da avaliacao foi observar, para diferentes perfis de entrada, como os custos
teoricos se refletem na pratica em termos de comparacoes, trocas e tempo.

Este documento foi revisado para atender as alteracoes das instrucoes de confeccao do
artigo, mantendo rastreabilidade entre texto, tabelas e arquivos de resultados.

\section{Aderencia ao Enunciado Atualizado}
O artigo atende aos pontos centrais definidos no enunciado:
\begin{itemize}
\item implementacao e avaliacao de 13 algoritmos de ordenacao;
\item comparacao em entradas de tamanho 10.000 e 1.000.000;
\item uso dos tres perfis de entrada (aleatoria, crescente e decrescente);
\item apresentacao das tres metricas exigidas: comparacoes, trocas e tempo;
\item inclusao das duas variantes de Quicksort exigidas (quicksortcentro e quicksortmediana);
\item uso de limite de 24 horas por algoritmo/instancia para classificar timeout.
\end{itemize}

\section{Metodologia Experimental}
\subsection{Ambiente e Instrumentacao}
Os experimentos foram executados em ambiente Windows com compilacao C11.
O programa \texttt{experimentos.c} foi utilizado para: (i) carregar as instancias,
(ii) executar o algoritmo selecionado, (iii) contabilizar comparacoes e trocas,
e (iv) medir tempo em segundos com \texttt{timespec\_get(TIME\_UTC)}.

Conforme o enunciado atualizado, foi adotado tempo maximo de 24 horas para cada
par algoritmo/instancia. Execucoes nao finalizadas nesse intervalo foram registradas
como \texttt{timeout} nos arquivos de resultados e nas tabelas deste artigo.

\subsection{Instancias e Cenarios}
Foram utilizadas seis instancias:
\begin{itemize}
\item 10.000 elementos aleatorios
\item 10.000 elementos ordenados crescentemente
\item 10.000 elementos ordenados decrescentemente
\item 1.000.000 elementos aleatorios
\item 1.000.000 elementos ordenados crescentemente
\item 1.000.000 elementos ordenados decrescentemente
\end{itemize}

Total de cenarios executados: $13 \times 6 = 78$.

\section{Resultados}
\subsection{Graficos de Tempo}
A Figura~\ref{fig:tempo10k} apresenta os tempos para $n=10.000$.
A Figura~\ref{fig:tempo1m} apresenta os tempos para $n=1.000.000$.
Valores \texttt{timeout} nao sao plotados (representados como \texttt{nan} nos CSVs
de graficos).

\begin{figure}[ht]
\centering
\begin{tikzpicture}
\begin{axis}[
width=\textwidth,
height=6.5cm,
ybar,
bar width=3pt,
ymajorgrids=true,
grid style=dashed,
unbounded coords=jump,
legend style={at={(0.5,1.15)},anchor=south,legend columns=3},
symbolic x coords={bolha,bolha-parada,insercao,insercao-bin,insercao-ter,shell,selecao,heap,quick-centro,quick-mediana,merge,radix,bucket},
xtick=data,
xticklabel style={rotate=45,anchor=east,font=\scriptsize},
ylabel={Tempo (s)},
xlabel={Algoritmo}
]
\addplot table [x=algoritmo,y=aleatoria,col sep=comma] {graficos/tempo_10000.csv};
\addplot table [x=algoritmo,y=crescente,col sep=comma] {graficos/tempo_10000.csv};
\addplot table [x=algoritmo,y=decrescente,col sep=comma] {graficos/tempo_10000.csv};
\legend{Aleatoria,Crescente,Decrescente}
\end{axis}
\end{tikzpicture}
\caption{Tempo de execucao para $n=10.000$.}
\label{fig:tempo10k}
\end{figure}

\begin{figure}[ht]
\centering
\begin{tikzpicture}
\begin{axis}[
width=\textwidth,
height=6.8cm,
ybar,
bar width=3pt,
ymode=log,
log basis y={10},
ymajorgrids=true,
grid style=dashed,
unbounded coords=jump,
legend style={at={(0.5,1.15)},anchor=south,legend columns=3},
symbolic x coords={bolha,bolha-parada,insercao,insercao-bin,insercao-ter,shell,selecao,heap,quick-centro,quick-mediana,merge,radix,bucket},
xtick=data,
xticklabel style={rotate=45,anchor=east,font=\scriptsize},
ylabel={Tempo (s) em escala log},
xlabel={Algoritmo}
]
\addplot table [x=algoritmo,y=aleatoria,col sep=comma] {graficos/tempo_1000000.csv};
\addplot table [x=algoritmo,y=crescente,col sep=comma] {graficos/tempo_1000000.csv};
\addplot table [x=algoritmo,y=decrescente,col sep=comma] {graficos/tempo_1000000.csv};
\legend{Aleatoria,Crescente,Decrescente}
\end{axis}
\end{tikzpicture}
\caption{Tempo de execucao para $n=1.000.000$ (escala logaritmica).}
\label{fig:tempo1m}
\end{figure}

\subsection{Graficos de Comparacoes e Trocas ($n=10.000$)}
As Figuras~\ref{fig:comp10k} e \ref{fig:trocas10k} destacam os custos internos
de comparacoes e trocas para $n=10.000$.

\begin{figure}[ht]
\centering
\begin{tikzpicture}
\begin{axis}[
width=\textwidth,
height=6.5cm,
ybar,
bar width=3pt,
ymode=log,
log basis y={10},
ymajorgrids=true,
grid style=dashed,
unbounded coords=jump,
legend style={at={(0.5,1.15)},anchor=south,legend columns=3},
symbolic x coords={bolha,bolha-parada,insercao,insercao-bin,insercao-ter,shell,selecao,heap,quick-centro,quick-mediana,merge,radix,bucket},
xtick=data,
xticklabel style={rotate=45,anchor=east,font=\scriptsize},
ylabel={Comparacoes (escala log)},
xlabel={Algoritmo}
]
\addplot table [x=algoritmo,y=aleatoria,col sep=comma] {graficos/comparacoes_10000.csv};
\addplot table [x=algoritmo,y=crescente,col sep=comma] {graficos/comparacoes_10000.csv};
\addplot table [x=algoritmo,y=decrescente,col sep=comma] {graficos/comparacoes_10000.csv};
\legend{Aleatoria,Crescente,Decrescente}
\end{axis}
\end{tikzpicture}
\caption{Numero de comparacoes para $n=10.000$.}
\label{fig:comp10k}
\end{figure}

\begin{figure}[ht]
\centering
\begin{tikzpicture}
\begin{axis}[
width=\textwidth,
height=6.5cm,
ybar,
bar width=3pt,
ymode=log,
log basis y={10},
ymajorgrids=true,
grid style=dashed,
unbounded coords=jump,
legend style={at={(0.5,1.15)},anchor=south,legend columns=3},
symbolic x coords={bolha,bolha-parada,insercao,insercao-bin,insercao-ter,shell,selecao,heap,quick-centro,quick-mediana,merge,radix,bucket},
xtick=data,
xticklabel style={rotate=45,anchor=east,font=\scriptsize},
ylabel={Trocas (escala log)},
xlabel={Algoritmo}
]
\addplot table [x=algoritmo,y=aleatoria,col sep=comma] {graficos/trocas_10000.csv};
\addplot table [x=algoritmo,y=crescente,col sep=comma] {graficos/trocas_10000.csv};
\addplot table [x=algoritmo,y=decrescente,col sep=comma] {graficos/trocas_10000.csv};
\legend{Aleatoria,Crescente,Decrescente}
\end{axis}
\end{tikzpicture}
\caption{Numero de trocas para $n=10.000$.}
\label{fig:trocas10k}
\end{figure}

\subsection{Tabelas Completas por Cenario}
As tabelas a seguir apresentam os resultados completos (comparacoes, trocas e tempo)
para cada tamanho e tipo de entrada.

\begin{table}[ht]
\centering
\caption{Resultados para $n=10.000$ com entrada aleatoria.}
\label{tab:10000a}
\begin{tabular}{lrrr}
\hline
Algoritmo & Comparacoes & Trocas & Tempo (s) \\
\hline
bolha & 49995000 & 25038974 & 0.026633024 \\
bolha-parada & 49980635 & 25038974 & 0.030628920 \\
insercao & 25048964 & 25048973 & 0.007726908 \\
insercao-bin & 118946 & 25048973 & 0.002248287 \\
insercao-ter & 127309 & 25048973 & 0.002520323 \\
shell & 263627 & 268749 & 0.000684977 \\
selecao & 49995000 & 9990 & 0.009216547 \\
heap & 235466 & 124254 & 0.000591516 \\
quick-centro & 172209 & 33856 & 0.000519037 \\
quick-mediana & 184058 & 35004 & 0.000505209 \\
merge & 120502 & 133616 & 0.001075268 \\
radix & 9999 & 100000 & 0.000152349 \\
bucket & 2540574 & 2550640 & 0.001029730 \\
\hline
\end{tabular}

\end{table}

\begin{table}[ht]
\centering
\caption{Resultados para $n=10.000$ com entrada crescente.}
\label{tab:10000c}
\begin{tabular}{lrrr}
\hline
Algoritmo & Comparacoes & Trocas & Tempo (s) \\
\hline
bolha & 49995000 & 0 & 0.009930372 \\
bolha-parada & 9999 & 0 & 0.000003099 \\
insercao & 9999 & 9999 & 0.000009537 \\
insercao-bin & 123617 & 9999 & 0.000074387 \\
insercao-ter & 150486 & 9999 & 0.000087023 \\
shell & 120005 & 120005 & 0.000089884 \\
selecao & 49995000 & 0 & 0.008524656 \\
heap & 244460 & 131956 & 0.000383615 \\
quick-centro & 125439 & 5904 & 0.000060081 \\
quick-mediana & 137247 & 5904 & 0.000067711 \\
merge & 69008 & 133616 & 0.000753164 \\
radix & 9999 & 100000 & 0.000179291 \\
bucket & 19989 & 29990 & 0.000050545 \\
\hline
\end{tabular}

\end{table}

\begin{table}[ht]
\centering
\caption{Resultados para $n=10.000$ com entrada decrescente.}
\label{tab:10000d}
\begin{tabular}{lrrr}
\hline
Algoritmo & Comparacoes & Trocas & Tempo (s) \\
\hline
bolha & 49995000 & 49995000 & 0.012869358 \\
bolha-parada & 49995000 & 49995000 & 0.014713287 \\
insercao & 49995000 & 50004999 & 0.015097141 \\
insercao-bin & 113631 & 50004999 & 0.003630400 \\
insercao-ter & 75243 & 50004999 & 0.003584385 \\
shell & 172578 & 182565 & 0.000130415 \\
selecao & 49995000 & 5000 & 0.009212971 \\
heap & 226682 & 116696 & 0.000426054 \\
quick-centro & 125452 & 10904 & 0.000068426 \\
quick-mediana & 137262 & 10904 & 0.000085592 \\
merge & 64608 & 133616 & 0.000753641 \\
radix & 9999 & 100000 & 0.000156879 \\
bucket & 5004999 & 5024990 & 0.001885891 \\
\hline
\end{tabular}

\end{table}

\begin{table}[ht]
\centering
\caption{Resultados para $n=1.000.000$ com entrada aleatoria.}
\label{tab:1000000a}
\begin{tabular}{lrrr}
\hline
Algoritmo & Comparacoes & Trocas & Tempo (s) \\
\hline
bolha & timeout & timeout & timeout \\
bolha-parada & timeout & timeout & timeout \\
insercao & timeout & timeout & timeout \\
insercao-bin & timeout & timeout & timeout \\
insercao-ter & timeout & timeout & timeout \\
shell & 65724076 & 66226717 & 0.126871824 \\
selecao & timeout & timeout & timeout \\
heap & 36794531 & 19049328 & 0.111006498 \\
quick-centro & 27903893 & 4852801 & 0.068959951 \\
quick-mediana & 27603809 & 5016553 & 0.069190264 \\
merge & 18674269 & 19951424 & 0.116347075 \\
radix & 999999 & 14000000 & 0.036098003 \\
bucket & timeout & timeout & timeout \\
\hline
\end{tabular}

\end{table}

\begin{table}[ht]
\centering
\caption{Resultados para $n=1.000.000$ com entrada crescente.}
\label{tab:1000000c}
\begin{tabular}{lrrr}
\hline
Algoritmo & Comparacoes & Trocas & Tempo (s) \\
\hline
bolha & timeout & timeout & timeout \\
bolha-parada & 999999 & 0 & 0.000366211 \\
insercao & 999999 & 999999 & 0.000937462 \\
insercao-bin & 18951425 & 999999 & 0.012098551 \\
insercao-ter & 23608530 & 999999 & 0.013160944 \\
shell & 18000007 & 18000007 & 0.013895988 \\
selecao & timeout & timeout & timeout \\
heap & 37692069 & 19787792 & 0.048102140 \\
quick-centro & 19000019 & 524287 & 0.007274866 \\
quick-mediana & 20048593 & 524287 & 0.007854223 \\
merge & 10066432 & 19951424 & 0.088405371 \\
radix & 999999 & 14000000 & 0.031106234 \\
bucket & 1999989 & 2999990 & 0.003863096 \\
\hline
\end{tabular}

\end{table}

\begin{table}[ht]
\centering
\caption{Resultados para $n=1.000.000$ com entrada decrescente.}
\label{tab:1000000d}
\begin{tabular}{lrrr}
\hline
Algoritmo & Comparacoes & Trocas & Tempo (s) \\
\hline
bolha & timeout & timeout & timeout \\
bolha-parada & timeout & timeout & timeout \\
insercao & timeout & timeout & timeout \\
insercao-bin & timeout & timeout & timeout \\
insercao-ter & timeout & timeout & timeout \\
shell & 26357530 & 27357511 & 0.017968178 \\
selecao & timeout & timeout & timeout \\
heap & 36001436 & 18333408 & 0.050829887 \\
quick-centro & 19000036 & 1024286 & 0.008491278 \\
quick-mediana & 20048610 & 1024286 & 0.008160114 \\
merge & 9884992 & 19951424 & 0.086262941 \\
radix & 999999 & 14000000 & 0.033104897 \\
bucket & timeout & timeout & timeout \\
\hline
\end{tabular}

\end{table}

\section{Discussao}
Os resultados confirmam o comportamento esperado dos algoritmos quadraticos.
Em $n=1.000.000$, bolha, selecao e insercoes gerais tornam-se inviaveis em
cenarios mais custosos, com ocorrencias de timeout registradas.

Nos casos onde a entrada ja esta ordenada, algoritmos adaptativos se beneficiam:
bolha com parada e insercao exibem tempos muito baixos em crescente.
Por outro lado, essa vantagem nao se sustenta em perfis desfavoraveis.

Entre os algoritmos de melhor escalabilidade, radix apresentou o menor tempo
na instancia aleatoria de 1.000.000 elementos (aproximadamente 0,036 s),
seguido de quicksort (centro e mediana) e heap/merge.

As duas variantes de quicksort tiveram desempenho semelhante nas instancias grandes.
Em entrada aleatoria, o pivo central foi ligeiramente mais rapido; em entrada
decrescente, a mediana de tres foi levemente melhor em tempo.
Esses dados indicam que ambas as abordagens sao competitivas, com pequenas variacoes
dependentes da distribuicao.

Bucket sort mostrou forte dependencia da distribuicao dos dados:
obteve bom resultado em crescente para 1.000.000, mas apresentou timeout em cenarios
aleatorio e decrescente nessa mesma escala.

\section{Conclusao}
O estudo mostrou, de forma quantitativa, como os custos assintoticos se refletem
na pratica sob diferentes entradas. Para volumes grandes, os metodos quadraticos
ficam restritos a casos muito favoraveis, enquanto quicksort, heap, merge e radix
mantem desempenho robusto.

Como continuidade, recomenda-se executar multiplas repeticoes por cenario,
reportar media e desvio-padrao e incluir analise estatistica para fortalecer
a comparabilidade entre implementacoes e a reproducibilidade experimental.

\end{document}

